% Generated draft: requires mathematical and source review.
\documentclass[11pt]{article}
\usepackage[margin=25mm]{geometry}
\usepackage{fontspec}
\setmainfont{Latin Modern Roman}
\usepackage{amsmath,amssymb,mathtools,unicode-math}
\setmathfont{Latin Modern Math}
\usepackage{enumitem,needspace,xurl,hyperref}
\setlength{\parindent}{0pt}
\setlength{\parskip}{5pt}
\setlength{\emergencystretch}{3em}
\newcommand{\N}{\mathbb{N}}
\newcommand{\Z}{\mathbb{Z}}
\newcommand{\Q}{\mathbb{Q}}
\newcommand{\R}{\mathbb{R}}
\newcommand{\C}{\mathbb{C}}
\newcommand{\F}{\mathbb{F}}
\newcommand{\A}{\mathbb{A}}
\newcommand{\PP}{\mathbb{P}}
\newcommand{\E}{\mathbb{E}}
\newcommand{\eps}{\varepsilon}
\newcommand{\dd}{\,\mathrm{d}}
\begin{document}
\section*{UnsolvedMath ID: 700467}
\textbf{Source problem code:} catalogue entry
\subsection{Definitions and mathematical statement}
The expressions \(\widetilde{N}\)-group and \(\overline{\! Z}\)-group denote the classes referenced in M. I. Kargapolov and Yu. I. Merzlyakov, \emph{Fundamentals of the Theory of Groups}, 3rd edition, Moscow, Nauka, 1982, p.~205 (Russian). The supplied record does not reproduce those definitions; their exact hypotheses and conventions therefore remain unresolved in this draft.

Subject to those referenced definitions, the question is whether
\[
\text{for every group }G,\qquad
\bigl(G\text{ is a }\widetilde{N}\text{-group}\bigr)
\Longrightarrow
\bigl(G\text{ is a }\overline{\! Z}\text{-group}\bigr).
\]
The record states no additional restriction on \(G\). This is a conditional transcription of the intended implication, requiring the cited definitions before it becomes a self-contained mathematical statement.
\subsection{Short English statement}
Does every \(\widetilde{N}\)-group belong to the class of \(\overline{\! Z}\)-groups, with both terms interpreted exactly as in the cited textbook? The definitions are absent from the supplied record and require review.
\subsection{Sources}
\begin{itemize}
\item Local UnsolvedMath registry, permanent ID \(700467\), source code \texttt{catalogue entry}, title \emph{Kourovka Notebook Problem 8.15}. The supplied registry labels the problem open; this label is not an independently verified assessment.
\item \emph{The Kourovka Notebook (Unsolved Problems in Group Theory)}, original problem 8.15, 8th issue (1982), as attributed in the record to the 21st issue edition (2026), \texttt{arXiv:1401.0300v48}. Proposed by Sh. S. Kemkhadze; the statement describes it as a well-known problem.
\item M. I. Kargapolov and Yu. I. Merzlyakov, \emph{Fundamentals of the Theory of Groups}, 3rd edition, Moscow, Nauka, 1982, p.~205 (Russian), cited by the supplied statement for the definitions of the two group classes.
\end{itemize}
\end{document}
